% The forward/backward entropy production, measured on masked evaluators.
% Draft: every number is a macro, generated by make_tables.py from the
% measurements, so nothing here is typed by hand.  Until those runs land the
% macros read [TODO:...] and make_arxiv.sh refuses to package.
%
% % The forward/backward entropy production, measured on masked evaluators.
% Draft: every number is a macro, generated by make_tables.py from the
% measurements, so nothing here is typed by hand.  Until those runs land the
% macros read [TODO:...] and make_arxiv.sh refuses to package.
%
% % The forward/backward entropy production, measured on masked evaluators.
% Draft: every number is a macro, generated by make_tables.py from the
% measurements, so nothing here is typed by hand.  Until those runs land the
% macros read [TODO:...] and make_arxiv.sh refuses to package.
%
% % The forward/backward entropy production, measured on masked evaluators.
% Draft: every number is a macro, generated by make_tables.py from the
% measurements, so nothing here is typed by hand.  Until those runs land the
% macros read [TODO:...] and make_arxiv.sh refuses to package.
%
% \input{sec_truesigma} after Section "Results II".

\section{Results III: $\sigma$, on masked evaluators}
\label{sec:truesigma}

Everything so far measures $\Delta$: a text scored against a rearrangement
of itself, which is the fixed-text reading of
Ref.~\cite{sagawa2026}~\cite[Sec.~V~B]{sagawa2026}.
Section~\ref{sec:notareversal} gave the reason it has no origin on the text
it is read from.  This section measures a comparison that has one.

\subsection{The estimator, and the identity that constrains it}
\label{sec:fbdef}

The estimator is Eq.~\eqref{eq:fb}: the any-order chain rule of
Eq.~\eqref{eq:aoar} evaluated at $\pi = $ left-to-right and again at $\pi
= $ right-to-left on the same text.  Equation~\eqref{eq:aoar} is a genuine
log-probability --- it normalises to one over sequences, for any fixed
$\pi$ --- and both of its evaluations score the same $x$.  What differs is the direction in which the
model assembles it: forward, one token at a time from the left, against
backward, one token at a time from the right.  That is the forward path
against its reverse, which is what the stochastic-thermodynamic definition
asks for.

Equation~\eqref{eq:fb} carries an identity that decides how it can be read,
stated as Eq.~\eqref{eq:fbzero} and proved here.  If the model's
conditionals are those of some joint distribution $P$, then every term of
Eq.~\eqref{eq:aoar} telescopes and $\log P_\pi(x) = \log P(x)$ for every
$\pi$; the two terms of Eq.~\eqref{eq:fb} are then equal, and $\sigma(x) =
0$ for every $x$, causal or not.  The size of a non-zero $\sigma$ is therefore
not a property of the text.  Where one arrangement of a text falls among
the rearrangements of that same text is a separate question, and
Section~\ref{sec:fbrank} is where it is asked.  Whatever
non-zero reading appears is a measure of the masked evaluator's \emph{in}compatibility
--- the failure of its conditionals to derive from any joint --- as exposed
by that text.  This is not an empirical finding to be checked; it is what
Eq.~\eqref{eq:aoar} says.

Two consequences follow immediately.

An autoregressive evaluator has nothing to report here, and it is worth
being exact about why, because the two available answers are different.  The
backward term of Eq.~\eqref{eq:fb} is $\log P_{\mathrm{r2l}}(x)$, and each of
its factors
conditions on a suffix while the prefix is still unknown, and a model that
attends only to what precedes emits none of them.  Equation~\eqref{eq:fb}
is well defined; what an autoregressive evaluator cannot do is supply the
conditionals it asks for.  So $\sigma$ is not evaluable on the
autoregressive evaluators Ref.~\cite{sagawa2026} uses.  What is absent is
not a number but an independent number: a masked model's backward term is
new information because its conditionals derive from no joint, and an
autoregressive model's conditionals derive from one.  That explains the
substitution of $\Delta$ for $\sigma$; it does not license reading $\Delta$
as $\sigma$.

A masked evaluator can be placed here, and is the only kind that can.  It
supplies every conditional Eq.~\eqref{eq:aoar} asks for, in any order.  That
is also the scope of Eq.~\eqref{eq:fbzero}: it says that $\sigma$ vanishes
for an evaluator whose conditionals derive from one joint, and an evaluator
that cannot supply those conditionals is not a case of it either way.  For a
masked evaluator no \emph{one} joint need account for all of the unmasking
orders --- that is what incompatibility is --- and where none does,
Eq.~\eqref{eq:fbzero} does not force $\sigma$ to zero.
Appendix~\ref{sec:toy} measured that incompatibility directly, as the
entropy production $\sigEP$ on a system small enough to enumerate, and
Table~\ref{tab:bayesroute} reads $\sigma$ itself on those same models:
$\sigEP$ and $\sigma$ are both non-zero at every width of that sweep, and
nothing in it relates their sizes (Section~\ref{sec:fbresult} gives the rank
correlation and the $n$ it rests on).  Here $\sigma$ is read off real text.
What the identity fixes is what a non-zero reading means --- that the
conditionals behind it derive from no single joint --- and not what the
value moves with.  Whether it also happens to separate the two label groups
is therefore not settled in advance, and has to be measured.

\subsection{Where the authored arrangement sits}
\label{sec:fbrank}

A single non-zero $\sigma(x)$ says that the conditionals behind it derive
from no single joint, and nothing more; it is not an amount of anything.
The question this section asks is narrower: whether the value $\sigma$
returns is nevertheless higher on the arrangement a text came in than on its
rearrangements, and for that the text has to be compared with itself.  Taking the same block arrangements as
Section~\ref{sec:blkord} --- the $n!$ orders of the same $n$ sentences --- and
writing $x_\pi$ for the text with its blocks in order $\pi$, we score
\begin{equation}
f(\pi) \;=\; \sigma(x_\pi)
\label{eq:fbperm}
\end{equation}
for every $\pi$, and record the midrank of $f(\mathrm{id})$, the authored arrangement,
among the $n!$ values.  Length, vocabulary, topic and style are identical
across the arrangements by construction, so the between-set confounds of
Section~\ref{sec:construct} cannot act.

The null is $0.50$, but only one of the two stages of Section~\ref{sec:blkord} carries over.  There the permutation set itself had
a centre, fixed at zero by the antisymmetry $\Delta(\pi) = -\Delta(\pi \circ
\mathrm{rev})$.  Equation~\eqref{eq:fbperm} has no such symmetry: reversing the arrangement changes which text is scored and does not negate $\sigma$,
so the set of $f(\pi)$ is not centred and supplies no origin of its own.  What survives is the second stage.  Under the hypothesis that the authored arrangement
is one arrangement among $n!$, its rank is uniform, so the midrank has
expectation $0.50$ whatever the shape of the $f$ distribution.

Two situations make the $n!$ values coincide, and both must be reported at
the null rather than ranked.  A masked evaluator whose conditionals derive
from one joint gives $f(\pi) = 0$ for every $\pi$ by
Eq.~\eqref{eq:fbzero}; in floating point those values differ by
$\sim\!10^{-16}$, and ranking them would manufacture a midrank out of
rounding noise.  A text whose sentences are identical strings has
arrangements that are literally the same token sequence.  Values within
$10^{-12}$ are therefore treated as tied, that text is flagged degenerate,
and its midrank is set to $0.50$ exactly.  Table~\ref{tab:fbrank} reports how
many scorings that was: \FbDegenTotal\ of \FbNScorings, over
\FbNTexts\ texts read by each of the two masked evaluators used here,
\texttt{bert-base-uncased} and \texttt{LLaDA-8B-Base}.

This is the \emph{within-text reading}: one text, its $n!$ arrangements
scored by Eq.~\eqref{eq:fbperm}, and the authored arrangement's midrank
among them.  Section~\ref{sec:fbresult} also reads $\sigma$ a second way,
without rearranging anything, and keeps the two apart throughout.


\subsection{What the measurement says}
\label{sec:fbresult}

\begin{table}[tbp]
\centering
{\small\input{tables/tab_fbrank}}
\caption{The forward/backward entropy production of
Eq.~\eqref{eq:fb}, read as the midrank of the authored arrangement among
the $n!$ arrangements of the same text's own sentences.  Columns from the
left: the masked evaluator; the corpus; $n$, the texts in that cell; how
many were degenerate (all $n!$ values tied, hence reported at the null and
excluded from $r$); the median midrank of the causal group and of the
non-causal group; $p_{c}$ and $p_{nc}$, the two-sided Wilcoxon signed-rank
$p$ of each group against the $0.50$ null; $r$, the rank-biserial
correlation of Eq.~\eqref{eq:rb} between the two groups, negative when the
causal group ranks lower; and its two-sided Mann--Whitney $p$.  The two
signed-rank columns answer a different question from the last one: whether
each group is displaced from the null at all, as against whether the groups
differ from each other.  Masked evaluators only
(Section~\ref{sec:truesigma}).}
\label{tab:fbrank}
\end{table}

Table~\ref{tab:fbrank} gives, for each masked evaluator and corpus, the median
midrank of the authored arrangement in each label group, the two-sided
Wilcoxon signed-rank $p$ of each group against the $0.50$ null in the
$p_{c}$ and $p_{nc}$ columns, and the effect size between the groups.  Two
readings come out of it, and they are of different kinds: where the authored
arrangement sits, which turns out to be a property of the masked evaluator
rather than of the text, and what the causal label adds to that, which turns
out to be almost nothing.

\begin{table}[tbp]
\centering
\input{tables/tab_fbdecomp}
\caption{The same measurement as Table~\ref{tab:fbrank}, summarised per
masked evaluator.  Columns from the left: the masked evaluator; the authored
component $\mathrm{med}(\mathrm{pct}_{\text{non-causal}}) - 0.50$, as a
median over the \NCorp\ corpora; that component in steps of the
$1/\GridK = \GridStep$ grid; the causal component of Eq.~\eqref{eq:caus},
likewise a median over corpora; that component in grid steps; and $r$, the
median over the \NCorp\ corpora of the rank-biserial correlation
Table~\ref{tab:fbrank} prints cell by cell, negative when the causal group
ranks lower.  These are the
numbers the text quotes; Table~\ref{tab:fbrank} is where they come from, one
corpus at a time.}
\label{tab:fbdecomp}
\end{table}

\paragraph{The authored arrangement's midrank is a property of the
masked evaluator.}
Equation~\eqref{eq:auth} is defined on midranks and says nothing about
which quantity produced them, so it reads under $\sigma$ exactly as it does
under $\Delta$.  Under \texttt{LLaDA} the authored component so read is
$\FbAuthoredLLaDA$: the authored arrangement sits that
far below the null, or $\FbAuthoredStepsLLaDA$ steps of the
$1/\GridK = \GridStep$ grid (Table~\ref{tab:fbdecomp}).  It is far from the null in both
label groups: \FbDispNRejectingLLaDA\ of its \FbNTestsDispLLaDA\
distance-from-null tests clear a Bonferroni threshold of
$\FbDispBonferroni$, set for the $\FbNTestsDisp$ such tests this section
reports, the largest of them at $p = \FbDispPMaxLLaDA$.  Those tests are
two-sided, so they establish that it is not at the null and not which side of it; the
direction is the sign of $\FbAuthoredLLaDA$ itself.  Under
\texttt{bert-base} the same texts give $\FbAuthoredBertBase$, which is
$\FbAuthoredStepsBertBase$ steps, and its distance-from-null tests clear that
threshold in $\FbDispNRejectingBertBase$ of $\FbNTestsDispBertBase$ cases.

The two readings stand $\FbStepGap$ grid steps apart on identical input.
Nothing about the texts changed between them; only the masked evaluator
changed.  A
quantity whose reading moves that far when the masked evaluator is replaced
is reporting a property of that evaluator.

The two masked evaluators do agree on one thing, and it is worth separating from
the disagreement.  The causally labelled group sits below the null in
$\FbCausalBelowBertBase$ of $\FbNCorporaOfBertBase$ corpora under
\texttt{bert-base} and in $\FbCausalBelowLLaDA$ of $\FbNCorporaOfLLaDA$
under \texttt{LLaDA}: where the causal group sits \emph{relative to the
null} is the same under both.  That is not a statement about the two groups'
order, which does reverse in one cell: \texttt{LLaDA} on \textsc{Opus46},
the only cell of Table~\ref{tab:fbrank} where $r$ is positive
($r = \FbRMax$).  The
\emph{level} both sit at is not, and the same two counts taken on the
non-causal group show it --- that group is above the null in
$\FbAuthoredAboveBertBase$ corpora under \texttt{bert-base} and in
$\FbAuthoredAboveLLaDA$ under \texttt{LLaDA}.  It is the level that moves by
$\FbStepGap$ steps.

\paragraph{The causal label does essentially nothing.}
Of the \FbNTestsCausal\ masked-evaluator--corpus tests of causal against non-causal,
\FbNSurviving\ survives a Bonferroni threshold of \FbBonferroni;
\FbNNominal\ fall below $0.05$ on their own, against \FbNExpected\ expected
by chance.  The sign is nonetheless the same in \FbNNegative\ of
\FbNCells\ cells, and we report it rather than only the failures: within
each masked evaluator, \texttt{LLaDA} gives \FbSignNegLLaDA\ negative ($p =
\FbSignPLLaDA$, exact binomial) and \texttt{bert-base}
\FbSignNegBertBase\ negative ($p = \FbSignPBertBase$).  Neither reaches
$0.05$ and neither could.  Under the null each corpus is an independent coin,
so the most extreme outcome --- all $\FbNCorporaOfLLaDA$ on one side --- has
probability $1/2^{\FbNCorporaOfLLaDA}$, and the two-sided $p$ of that
outcome is $2/2^{\FbNCorporaOfLLaDA} = \FbSignPBertBase$: the smallest the
test can return at this $n$, and the value \texttt{bert-base} takes.  Pooling the
\FbNCells\ cells would be the wrong test: both masked evaluators score the
\emph{same} texts, so the cells are not independent, and the corpus is the
unit that is.  The one surviving cell is \texttt{bert-base} on
\textsc{Gemini31} --- which is also the corpus where \texttt{LLaDA}'s distance from the null is largest, so it is more plausibly a property of
that corpus than of the label.

Taken together: the authored component under $\sigma$ moves
$\FbStepGap$ grid steps when the masked evaluator is changed, while
exchanging the causal and non-causal labels leaves \FbNSurviving\ of
\FbNCells\ cells surviving correction.  The separation between causally and
arbitrarily ordered texts that $\Delta$ shows in \ContrastNRejecting\ of
\ContrastNCells\ cells does not reappear here.  This is
Eq.~\eqref{eq:fbzero} seen from the measurement side.  The identity says
that where the conditionals derive from a joint there is no value range at
all; what the measurement adds is that where they do not --- which is the
case for \texttt{bert-base} and \texttt{LLaDA} alike --- what moves inside
the range that opens is the evaluator and not the label.

\begin{table}[tbp]
\centering
\input{tables/tab_fbbetween}
\caption{The unrearranged reading of Eq.~\eqref{eq:fb}: each text scored
once in each direction, in the arrangement it came in, with no permutation
set and no midrank.  \emph{Upper block}: the median $\sigma$ over all of a
cell's texts, both label groups together, per token.  It is not the authored
component of Eq.~\eqref{eq:auth}, which is a midrank against a null of
$0.50$; the null here is $\sigma = 0$.  \emph{Lower block}: the causal
component of this reading --- the median $\sigma$ of the causally labelled
texts less that of the non-causal ones --- which shares its name with
Eq.~\eqref{eq:caus} but is a difference of raw $\sigma$ values rather than
of midranks.  For a masked evaluator whose conditionals derive from a joint
every entry of both blocks would be exactly zero by Eq.~\eqref{eq:fbzero}.
None is, and on \FbBetweenNOpposite\ of the \FbBetweenNCompared\ corpora
the two masked evaluators give the upper block's median $\sigma$ opposite
signs.}
\label{tab:fbbetween}
\end{table}

A second reading makes the same point more directly.  Call it the
\emph{unrearranged reading}: Eq.~\eqref{eq:fb} evaluated as it stands, on
the text in the arrangement it came in, with no permutation set behind it.
Across cells the median $\sigma$
runs $\FbBetweenMinLLaDA$ to $\FbBetweenMaxLLaDA$ per token under
\texttt{LLaDA} and $\FbBetweenMinBertBase$ to $\FbBetweenMaxBertBase$ under
\texttt{bert-base}, positive in \FbBetweenNPosLLaDA\ of
\FbBetweenNOfLLaDA\ corpora under the first and \FbBetweenNPosBertBase\ of
\FbBetweenNOfBertBase\ under the second (Table~\ref{tab:fbbetween}).  On
\FbBetweenNOpposite\ of the \FbBetweenNCompared\ corpora the two give that
median opposite signs.  The texts are identical; what differs is which masked
evaluator computed $\sigma$.  What Eq.~\eqref{eq:fbzero} settles before any
corpus is scored is that a non-zero reading belongs to the evaluator's
incompatibility rather than to the text; a disagreement in sign on identical
texts is one way that shows, and confirms it rather than establishing it.  The causal
component of the same reading is the lower block of
Table~\ref{tab:fbbetween}: a median of $\FbBetweenCausalLLaDA$ under
\texttt{LLaDA} and $\FbBetweenCausalBertBase$ under \texttt{bert-base},
negative in \FbBetweenNNegOfLLaDA\ and \FbBetweenNNegOfBertBase\ corpora
respectively.  Both readings therefore give the causal group's median less
the non-causal group's the same sign, as a median over corpora: the causally
labelled texts sit lower.  The two are not in the same units: the
within-text reading is a difference of midranks ($\FbCausalLLaDA$ and
$\FbCausalBertBase$, or $\FbCausalStepsLLaDA$ and
$\FbCausalStepsBertBase$ grid steps), the unrearranged reading a difference
of raw $\sigma$.  The agreement is therefore one of direction and not of
size, and neither reaches significance.

\paragraph{Why this is not the argument made against $\Delta$.}
The two negative results in this paper are not of the same kind, and the
distinction decides what either of them licenses.  Under $\Delta$
(Sections~\ref{sec:authored} and~\ref{sec:causal}) the two label groups do
separate, in \ContrastNRejecting\ of \ContrastNCells\ cells.  What cannot be
read off is the size of the part attributable to the causal label: it sits at
the grid's resolution, and the groups have no criterion distinguishing them
beyond the prompts that produced them.  Under $\sigma$
(this section) what a non-zero reading reports is the masked evaluator's
incompatibility, by the identity, and not causal order.  The $\Delta$
result is empirical; the $\sigma$ one follows from the identity.  What $\sigma$ reports is
a defect of the masked evaluator: Eq.~\eqref{eq:fbzero} makes it identically zero
wherever the conditionals derive from a joint, so a reading of any size is
a reading of incompatibility, and a contrast between text sets under
$\sigma$ would be a contrast in how much of that incompatibility each set
exposes.  This is an argument about the quantity and it holds before any
measurement.  The direction differs between the two: under $\Delta$ the
causal group scores higher, under $\sigma$ lower.  We do not account for
that difference, and no claim in this paper needs to --- by Eq.~\eqref{eq:fbzero} a
reading of $\sigma$ of any size, in either direction, is a reading of the
masked evaluator's incompatibility and not of the text's causal structure.

Nor does the sweep of Appendix~\ref{sec:toy} show $\sigma$ heading towards
zero as the transformer width $d$ grows from $\CtrlDMin$ to $\CtrlDMax$.
The appendix reports two quantities on that one sweep, $\sigEP$ and
$\sigma$, which is why both are in view here; the statement that bears on
this section has to be the one about $\sigma$.  They are separate
quantities, each computed from the evaluator's own conditionals, and what
makes each of them non-zero is the same incompatibility.  The behaviour of
one does not follow from the behaviour of the other, and the sweep does not
show them moving together: ranked against each other they give
$\rho = \CtrlRhoXS$ at $p = \CtrlPXS$, which establishes nothing either way
at $n = \CtrlNWidths$.  What $\sigma$ does over the sweep is
not approach zero: $\langle|\sigma|\rangle$ runs between $\CtrlSigAbsMin$
and $\CtrlSigAbsMax$ across all $\CtrlNWidths$ widths.  By the
standard of Section~\ref{sec:window} that is as far as it goes --- a sweep
of $\CtrlNWidths$ points cannot establish that a quantity is constant, and
we do not claim the incompatibility is permanent.  What we claim is that
nothing measured here is approaching the compatible limit.

None of this transfers to $\Delta$.  On a fixed text there is no limit in
which $\Delta$ becomes zero and no identity to appeal to: it is a
log-ratio between a text
and its reversal, and nothing makes that ratio vanish as the autoregressive
evaluator's parameter count grows along the AR ladder.  Measured, its magnitude grows rather than shrinks along the AR ladder
(Section~\ref{sec:capacity}).  The grounds for declining to read $\Delta$ as
a measurement of causal order are the ones given there, in
Section~\ref{sec:resolution} and in Section~\ref{sec:construct}, and they
are empirical: they would stand even if everything this section says about
$\sigma$ were wrong.

The within-text reading and the unrearranged reading are computed by
separate programs, \fn{block_sigma_fb.py} and
\fn{order_intervention.py}.  They are not independent implementations: the
first imports the text loader and the masked-model provider from the second,
so what they share is the input path and the model interface, and what
differs is the construction of $\sigma$ itself.  On the authored
arrangement the median $\sigma$ they return agrees to the printed precision
in all \FbNCorpora\ corpora, which is the check that the rank construction
measures what the direct one does (\fn{check_two_constructions.py}).  That
check covers \texttt{LLaDA} only: no direct run exists for
\texttt{bert-base}.

\subsection{The mask rate \texttt{bert-base} was trained at}
\label{sec:maskrate}

One objection applies to \texttt{bert-base} and not to \texttt{LLaDA}.
Equation~\eqref{eq:aoar} sweeps the fraction of masked positions from $1$
down to $1/L$ as the chain rule reveals tokens.  LLaDA is trained with the
masking fraction drawn uniformly, so every term is in distribution.
\textsc{bert} is trained at a single rate, $0.15$, so most terms of
Eq.~\eqref{eq:aoar} ask it for a conditional at a rate it never saw.  A
natural worry is that the chain rule simply collapses there, and that
anything computed from it is noise.

It does not collapse.  Across all \FbNTexts\ texts and both masked evaluators ---
\FbNScorings\ scorings --- \FbDegenTotal\ produced $n!$ values that were
all equal.  That is what the column detects and all that it detects: a text
whose arrangements the masked evaluator cannot separate at all would have no
midrank, and none occurred, at any mask rate asked for.  It is not a claim
that the $n!$ values are distinct from one another, which this test does not
examine.  The reading is also not degenerate in the weaker
sense: within texts \texttt{bert-base}'s causally labelled group
sits below the null in $\FbCausalBelowBertBase$ of
$\FbNCorporaOfBertBase$ corpora, as \texttt{LLaDA}'s does in
$\FbCausalBelowLLaDA$ of $\FbNCorporaOfLLaDA$.  What differs between the two
masked evaluators within texts is the level the groups sit at, not the order
they sit in.

How far the authored arrangement sits from the null under $\sigma$ is the
quantity Section~\ref{sec:fbresult} measured, and what it found was that
its value tracks the masked evaluator and not the label.  We do not claim
\textsc{bert}'s out-of-distribution conditionals are accurate.  We claim only
that they are usable, that the collapse the objection anticipates does not
occur, and that the conclusion of Section~\ref{sec:fbdef} rests on an
identity rather than on either masked evaluator.
 after Section "Results II".

\section{Results III: $\sigma$, on masked evaluators}
\label{sec:truesigma}

Everything so far measures $\Delta$: a text scored against a rearrangement
of itself, which is the fixed-text reading of
Ref.~\cite{sagawa2026}~\cite[Sec.~V~B]{sagawa2026}.
Section~\ref{sec:notareversal} gave the reason it has no origin on the text
it is read from.  This section measures a comparison that has one.

\subsection{The estimator, and the identity that constrains it}
\label{sec:fbdef}

The estimator is Eq.~\eqref{eq:fb}: the any-order chain rule of
Eq.~\eqref{eq:aoar} evaluated at $\pi = $ left-to-right and again at $\pi
= $ right-to-left on the same text.  Equation~\eqref{eq:aoar} is a genuine
log-probability --- it normalises to one over sequences, for any fixed
$\pi$ --- and both of its evaluations score the same $x$.  What differs is the direction in which the
model assembles it: forward, one token at a time from the left, against
backward, one token at a time from the right.  That is the forward path
against its reverse, which is what the stochastic-thermodynamic definition
asks for.

Equation~\eqref{eq:fb} carries an identity that decides how it can be read,
stated as Eq.~\eqref{eq:fbzero} and proved here.  If the model's
conditionals are those of some joint distribution $P$, then every term of
Eq.~\eqref{eq:aoar} telescopes and $\log P_\pi(x) = \log P(x)$ for every
$\pi$; the two terms of Eq.~\eqref{eq:fb} are then equal, and $\sigma(x) =
0$ for every $x$, causal or not.  The size of a non-zero $\sigma$ is therefore
not a property of the text.  Where one arrangement of a text falls among
the rearrangements of that same text is a separate question, and
Section~\ref{sec:fbrank} is where it is asked.  Whatever
non-zero reading appears is a measure of the masked evaluator's \emph{in}compatibility
--- the failure of its conditionals to derive from any joint --- as exposed
by that text.  This is not an empirical finding to be checked; it is what
Eq.~\eqref{eq:aoar} says.

Two consequences follow immediately.

An autoregressive evaluator has nothing to report here, and it is worth
being exact about why, because the two available answers are different.  The
backward term of Eq.~\eqref{eq:fb} is $\log P_{\mathrm{r2l}}(x)$, and each of
its factors
conditions on a suffix while the prefix is still unknown, and a model that
attends only to what precedes emits none of them.  Equation~\eqref{eq:fb}
is well defined; what an autoregressive evaluator cannot do is supply the
conditionals it asks for.  So $\sigma$ is not evaluable on the
autoregressive evaluators Ref.~\cite{sagawa2026} uses.  What is absent is
not a number but an independent number: a masked model's backward term is
new information because its conditionals derive from no joint, and an
autoregressive model's conditionals derive from one.  That explains the
substitution of $\Delta$ for $\sigma$; it does not license reading $\Delta$
as $\sigma$.

A masked evaluator can be placed here, and is the only kind that can.  It
supplies every conditional Eq.~\eqref{eq:aoar} asks for, in any order.  That
is also the scope of Eq.~\eqref{eq:fbzero}: it says that $\sigma$ vanishes
for an evaluator whose conditionals derive from one joint, and an evaluator
that cannot supply those conditionals is not a case of it either way.  For a
masked evaluator no \emph{one} joint need account for all of the unmasking
orders --- that is what incompatibility is --- and where none does,
Eq.~\eqref{eq:fbzero} does not force $\sigma$ to zero.
Appendix~\ref{sec:toy} measured that incompatibility directly, as the
entropy production $\sigEP$ on a system small enough to enumerate, and
Table~\ref{tab:bayesroute} reads $\sigma$ itself on those same models:
$\sigEP$ and $\sigma$ are both non-zero at every width of that sweep, and
nothing in it relates their sizes (Section~\ref{sec:fbresult} gives the rank
correlation and the $n$ it rests on).  Here $\sigma$ is read off real text.
What the identity fixes is what a non-zero reading means --- that the
conditionals behind it derive from no single joint --- and not what the
value moves with.  Whether it also happens to separate the two label groups
is therefore not settled in advance, and has to be measured.

\subsection{Where the authored arrangement sits}
\label{sec:fbrank}

A single non-zero $\sigma(x)$ says that the conditionals behind it derive
from no single joint, and nothing more; it is not an amount of anything.
The question this section asks is narrower: whether the value $\sigma$
returns is nevertheless higher on the arrangement a text came in than on its
rearrangements, and for that the text has to be compared with itself.  Taking the same block arrangements as
Section~\ref{sec:blkord} --- the $n!$ orders of the same $n$ sentences --- and
writing $x_\pi$ for the text with its blocks in order $\pi$, we score
\begin{equation}
f(\pi) \;=\; \sigma(x_\pi)
\label{eq:fbperm}
\end{equation}
for every $\pi$, and record the midrank of $f(\mathrm{id})$, the authored arrangement,
among the $n!$ values.  Length, vocabulary, topic and style are identical
across the arrangements by construction, so the between-set confounds of
Section~\ref{sec:construct} cannot act.

The null is $0.50$, but only one of the two stages of Section~\ref{sec:blkord} carries over.  There the permutation set itself had
a centre, fixed at zero by the antisymmetry $\Delta(\pi) = -\Delta(\pi \circ
\mathrm{rev})$.  Equation~\eqref{eq:fbperm} has no such symmetry: reversing the arrangement changes which text is scored and does not negate $\sigma$,
so the set of $f(\pi)$ is not centred and supplies no origin of its own.  What survives is the second stage.  Under the hypothesis that the authored arrangement
is one arrangement among $n!$, its rank is uniform, so the midrank has
expectation $0.50$ whatever the shape of the $f$ distribution.

Two situations make the $n!$ values coincide, and both must be reported at
the null rather than ranked.  A masked evaluator whose conditionals derive
from one joint gives $f(\pi) = 0$ for every $\pi$ by
Eq.~\eqref{eq:fbzero}; in floating point those values differ by
$\sim\!10^{-16}$, and ranking them would manufacture a midrank out of
rounding noise.  A text whose sentences are identical strings has
arrangements that are literally the same token sequence.  Values within
$10^{-12}$ are therefore treated as tied, that text is flagged degenerate,
and its midrank is set to $0.50$ exactly.  Table~\ref{tab:fbrank} reports how
many scorings that was: \FbDegenTotal\ of \FbNScorings, over
\FbNTexts\ texts read by each of the two masked evaluators used here,
\texttt{bert-base-uncased} and \texttt{LLaDA-8B-Base}.

This is the \emph{within-text reading}: one text, its $n!$ arrangements
scored by Eq.~\eqref{eq:fbperm}, and the authored arrangement's midrank
among them.  Section~\ref{sec:fbresult} also reads $\sigma$ a second way,
without rearranging anything, and keeps the two apart throughout.


\subsection{What the measurement says}
\label{sec:fbresult}

\begin{table}[tbp]
\centering
{\small% generated by make_tables.py -- do not edit
\begin{tabular}{llrrrrrrrr}
\toprule
evaluator & corpus & $n$ & degen & causal & non-causal & $p_{c}$ & $p_{nc}$ & $r$ & $p$ \\
\midrule
bert-base-uncased & Opus46 & 200 & 0 & $0.458$ & $0.521$ & $0.366$ & $0.968$ & $-0.060$ & $0.462$ \\
bert-base-uncased & GPT54 & 200 & 0 & $0.479$ & $0.521$ & $0.676$ & $0.514$ & $-0.061$ & $0.454$ \\
bert-base-uncased & GPT56 & 200 & 0 & $0.479$ & $0.583$ & $0.989$ & $0.037$ & $-0.126$ & $0.122$ \\
bert-base-uncased & Gemini31 & 200 & 0 & $0.292$ & $0.479$ & $6\times10^{-6}$ & $0.759$ & $-0.287$ & $4\times10^{-4}$ \\
bert-base-uncased & Fable5 & 200 & 0 & $0.417$ & $0.521$ & $0.050$ & $0.736$ & $-0.130$ & $0.112$ \\
LLaDA-8B-Base & Opus46 & 200 & 0 & $0.292$ & $0.229$ & $6\times10^{-6}$ & $1\times10^{-11}$ & $+0.107$ & $0.192$ \\
LLaDA-8B-Base & GPT54 & 200 & 0 & $0.167$ & $0.188$ & $3\times10^{-15}$ & $9\times10^{-14}$ & $-0.187$ & $0.022$ \\
LLaDA-8B-Base & GPT56 & 200 & 0 & $0.229$ & $0.271$ & $2\times10^{-11}$ & $4\times10^{-9}$ & $-0.111$ & $0.175$ \\
LLaDA-8B-Base & Gemini31 & 200 & 0 & $0.146$ & $0.188$ & $9\times10^{-15}$ & $1\times10^{-13}$ & $-0.128$ & $0.115$ \\
LLaDA-8B-Base & Fable5 & 200 & 0 & $0.208$ & $0.229$ & $7\times10^{-11}$ & $6\times10^{-12}$ & $-0.047$ & $0.563$ \\
\bottomrule
\end{tabular}
}
\caption{The forward/backward entropy production of
Eq.~\eqref{eq:fb}, read as the midrank of the authored arrangement among
the $n!$ arrangements of the same text's own sentences.  Columns from the
left: the masked evaluator; the corpus; $n$, the texts in that cell; how
many were degenerate (all $n!$ values tied, hence reported at the null and
excluded from $r$); the median midrank of the causal group and of the
non-causal group; $p_{c}$ and $p_{nc}$, the two-sided Wilcoxon signed-rank
$p$ of each group against the $0.50$ null; $r$, the rank-biserial
correlation of Eq.~\eqref{eq:rb} between the two groups, negative when the
causal group ranks lower; and its two-sided Mann--Whitney $p$.  The two
signed-rank columns answer a different question from the last one: whether
each group is displaced from the null at all, as against whether the groups
differ from each other.  Masked evaluators only
(Section~\ref{sec:truesigma}).}
\label{tab:fbrank}
\end{table}

Table~\ref{tab:fbrank} gives, for each masked evaluator and corpus, the median
midrank of the authored arrangement in each label group, the two-sided
Wilcoxon signed-rank $p$ of each group against the $0.50$ null in the
$p_{c}$ and $p_{nc}$ columns, and the effect size between the groups.  Two
readings come out of it, and they are of different kinds: where the authored
arrangement sits, which turns out to be a property of the masked evaluator
rather than of the text, and what the causal label adds to that, which turns
out to be almost nothing.

\begin{table}[tbp]
\centering
% generated by make_tables.py -- do not edit
\begin{tabular}{lrrrrr}
\toprule
evaluator & authored & steps & causal & steps & $r$ \\
\midrule
bert-base-uncased & $+0.0208$ & $+1.0$ & $-0.1042$ & $-5.0$ & $-0.126$ \\
LLaDA-8B-Base & $-0.2708$ & $-13.0$ & $-0.0208$ & $-1.0$ & $-0.111$ \\
\bottomrule
\end{tabular}

\caption{The same measurement as Table~\ref{tab:fbrank}, summarised per
masked evaluator.  Columns from the left: the masked evaluator; the authored
component $\mathrm{med}(\mathrm{pct}_{\text{non-causal}}) - 0.50$, as a
median over the \NCorp\ corpora; that component in steps of the
$1/\GridK = \GridStep$ grid; the causal component of Eq.~\eqref{eq:caus},
likewise a median over corpora; that component in grid steps; and $r$, the
median over the \NCorp\ corpora of the rank-biserial correlation
Table~\ref{tab:fbrank} prints cell by cell, negative when the causal group
ranks lower.  These are the
numbers the text quotes; Table~\ref{tab:fbrank} is where they come from, one
corpus at a time.}
\label{tab:fbdecomp}
\end{table}

\paragraph{The authored arrangement's midrank is a property of the
masked evaluator.}
Equation~\eqref{eq:auth} is defined on midranks and says nothing about
which quantity produced them, so it reads under $\sigma$ exactly as it does
under $\Delta$.  Under \texttt{LLaDA} the authored component so read is
$\FbAuthoredLLaDA$: the authored arrangement sits that
far below the null, or $\FbAuthoredStepsLLaDA$ steps of the
$1/\GridK = \GridStep$ grid (Table~\ref{tab:fbdecomp}).  It is far from the null in both
label groups: \FbDispNRejectingLLaDA\ of its \FbNTestsDispLLaDA\
distance-from-null tests clear a Bonferroni threshold of
$\FbDispBonferroni$, set for the $\FbNTestsDisp$ such tests this section
reports, the largest of them at $p = \FbDispPMaxLLaDA$.  Those tests are
two-sided, so they establish that it is not at the null and not which side of it; the
direction is the sign of $\FbAuthoredLLaDA$ itself.  Under
\texttt{bert-base} the same texts give $\FbAuthoredBertBase$, which is
$\FbAuthoredStepsBertBase$ steps, and its distance-from-null tests clear that
threshold in $\FbDispNRejectingBertBase$ of $\FbNTestsDispBertBase$ cases.

The two readings stand $\FbStepGap$ grid steps apart on identical input.
Nothing about the texts changed between them; only the masked evaluator
changed.  A
quantity whose reading moves that far when the masked evaluator is replaced
is reporting a property of that evaluator.

The two masked evaluators do agree on one thing, and it is worth separating from
the disagreement.  The causally labelled group sits below the null in
$\FbCausalBelowBertBase$ of $\FbNCorporaOfBertBase$ corpora under
\texttt{bert-base} and in $\FbCausalBelowLLaDA$ of $\FbNCorporaOfLLaDA$
under \texttt{LLaDA}: where the causal group sits \emph{relative to the
null} is the same under both.  That is not a statement about the two groups'
order, which does reverse in one cell: \texttt{LLaDA} on \textsc{Opus46},
the only cell of Table~\ref{tab:fbrank} where $r$ is positive
($r = \FbRMax$).  The
\emph{level} both sit at is not, and the same two counts taken on the
non-causal group show it --- that group is above the null in
$\FbAuthoredAboveBertBase$ corpora under \texttt{bert-base} and in
$\FbAuthoredAboveLLaDA$ under \texttt{LLaDA}.  It is the level that moves by
$\FbStepGap$ steps.

\paragraph{The causal label does essentially nothing.}
Of the \FbNTestsCausal\ masked-evaluator--corpus tests of causal against non-causal,
\FbNSurviving\ survives a Bonferroni threshold of \FbBonferroni;
\FbNNominal\ fall below $0.05$ on their own, against \FbNExpected\ expected
by chance.  The sign is nonetheless the same in \FbNNegative\ of
\FbNCells\ cells, and we report it rather than only the failures: within
each masked evaluator, \texttt{LLaDA} gives \FbSignNegLLaDA\ negative ($p =
\FbSignPLLaDA$, exact binomial) and \texttt{bert-base}
\FbSignNegBertBase\ negative ($p = \FbSignPBertBase$).  Neither reaches
$0.05$ and neither could.  Under the null each corpus is an independent coin,
so the most extreme outcome --- all $\FbNCorporaOfLLaDA$ on one side --- has
probability $1/2^{\FbNCorporaOfLLaDA}$, and the two-sided $p$ of that
outcome is $2/2^{\FbNCorporaOfLLaDA} = \FbSignPBertBase$: the smallest the
test can return at this $n$, and the value \texttt{bert-base} takes.  Pooling the
\FbNCells\ cells would be the wrong test: both masked evaluators score the
\emph{same} texts, so the cells are not independent, and the corpus is the
unit that is.  The one surviving cell is \texttt{bert-base} on
\textsc{Gemini31} --- which is also the corpus where \texttt{LLaDA}'s distance from the null is largest, so it is more plausibly a property of
that corpus than of the label.

Taken together: the authored component under $\sigma$ moves
$\FbStepGap$ grid steps when the masked evaluator is changed, while
exchanging the causal and non-causal labels leaves \FbNSurviving\ of
\FbNCells\ cells surviving correction.  The separation between causally and
arbitrarily ordered texts that $\Delta$ shows in \ContrastNRejecting\ of
\ContrastNCells\ cells does not reappear here.  This is
Eq.~\eqref{eq:fbzero} seen from the measurement side.  The identity says
that where the conditionals derive from a joint there is no value range at
all; what the measurement adds is that where they do not --- which is the
case for \texttt{bert-base} and \texttt{LLaDA} alike --- what moves inside
the range that opens is the evaluator and not the label.

\begin{table}[tbp]
\centering
% generated by make_tables.py -- do not edit
\begin{tabular}{lrr}
\toprule
corpus & bert-base-uncased & LLaDA-8B-Base \\
\midrule
Opus46 & $-0.1235$ & $+0.0574$ \\
GPT54 & $-0.0169$ & $+0.0753$ \\
GPT56 & $+0.0864$ & $+0.1088$ \\
Gemini31 & $-0.1191$ & $+0.0161$ \\
Fable5 & $-0.0384$ & $+0.0552$ \\
\midrule
\multicolumn{3}{l}{\emph{causal component}} \\
Opus46 & $-0.2712$ & $+0.0127$ \\
GPT54 & $-0.4284$ & $-0.0295$ \\
GPT56 & $-0.3251$ & $-0.0500$ \\
Gemini31 & $-0.1727$ & $-0.0284$ \\
Fable5 & $-0.3798$ & $-0.0552$ \\
\bottomrule
\end{tabular}

\caption{The unrearranged reading of Eq.~\eqref{eq:fb}: each text scored
once in each direction, in the arrangement it came in, with no permutation
set and no midrank.  \emph{Upper block}: the median $\sigma$ over all of a
cell's texts, both label groups together, per token.  It is not the authored
component of Eq.~\eqref{eq:auth}, which is a midrank against a null of
$0.50$; the null here is $\sigma = 0$.  \emph{Lower block}: the causal
component of this reading --- the median $\sigma$ of the causally labelled
texts less that of the non-causal ones --- which shares its name with
Eq.~\eqref{eq:caus} but is a difference of raw $\sigma$ values rather than
of midranks.  For a masked evaluator whose conditionals derive from a joint
every entry of both blocks would be exactly zero by Eq.~\eqref{eq:fbzero}.
None is, and on \FbBetweenNOpposite\ of the \FbBetweenNCompared\ corpora
the two masked evaluators give the upper block's median $\sigma$ opposite
signs.}
\label{tab:fbbetween}
\end{table}

A second reading makes the same point more directly.  Call it the
\emph{unrearranged reading}: Eq.~\eqref{eq:fb} evaluated as it stands, on
the text in the arrangement it came in, with no permutation set behind it.
Across cells the median $\sigma$
runs $\FbBetweenMinLLaDA$ to $\FbBetweenMaxLLaDA$ per token under
\texttt{LLaDA} and $\FbBetweenMinBertBase$ to $\FbBetweenMaxBertBase$ under
\texttt{bert-base}, positive in \FbBetweenNPosLLaDA\ of
\FbBetweenNOfLLaDA\ corpora under the first and \FbBetweenNPosBertBase\ of
\FbBetweenNOfBertBase\ under the second (Table~\ref{tab:fbbetween}).  On
\FbBetweenNOpposite\ of the \FbBetweenNCompared\ corpora the two give that
median opposite signs.  The texts are identical; what differs is which masked
evaluator computed $\sigma$.  What Eq.~\eqref{eq:fbzero} settles before any
corpus is scored is that a non-zero reading belongs to the evaluator's
incompatibility rather than to the text; a disagreement in sign on identical
texts is one way that shows, and confirms it rather than establishing it.  The causal
component of the same reading is the lower block of
Table~\ref{tab:fbbetween}: a median of $\FbBetweenCausalLLaDA$ under
\texttt{LLaDA} and $\FbBetweenCausalBertBase$ under \texttt{bert-base},
negative in \FbBetweenNNegOfLLaDA\ and \FbBetweenNNegOfBertBase\ corpora
respectively.  Both readings therefore give the causal group's median less
the non-causal group's the same sign, as a median over corpora: the causally
labelled texts sit lower.  The two are not in the same units: the
within-text reading is a difference of midranks ($\FbCausalLLaDA$ and
$\FbCausalBertBase$, or $\FbCausalStepsLLaDA$ and
$\FbCausalStepsBertBase$ grid steps), the unrearranged reading a difference
of raw $\sigma$.  The agreement is therefore one of direction and not of
size, and neither reaches significance.

\paragraph{Why this is not the argument made against $\Delta$.}
The two negative results in this paper are not of the same kind, and the
distinction decides what either of them licenses.  Under $\Delta$
(Sections~\ref{sec:authored} and~\ref{sec:causal}) the two label groups do
separate, in \ContrastNRejecting\ of \ContrastNCells\ cells.  What cannot be
read off is the size of the part attributable to the causal label: it sits at
the grid's resolution, and the groups have no criterion distinguishing them
beyond the prompts that produced them.  Under $\sigma$
(this section) what a non-zero reading reports is the masked evaluator's
incompatibility, by the identity, and not causal order.  The $\Delta$
result is empirical; the $\sigma$ one follows from the identity.  What $\sigma$ reports is
a defect of the masked evaluator: Eq.~\eqref{eq:fbzero} makes it identically zero
wherever the conditionals derive from a joint, so a reading of any size is
a reading of incompatibility, and a contrast between text sets under
$\sigma$ would be a contrast in how much of that incompatibility each set
exposes.  This is an argument about the quantity and it holds before any
measurement.  The direction differs between the two: under $\Delta$ the
causal group scores higher, under $\sigma$ lower.  We do not account for
that difference, and no claim in this paper needs to --- by Eq.~\eqref{eq:fbzero} a
reading of $\sigma$ of any size, in either direction, is a reading of the
masked evaluator's incompatibility and not of the text's causal structure.

Nor does the sweep of Appendix~\ref{sec:toy} show $\sigma$ heading towards
zero as the transformer width $d$ grows from $\CtrlDMin$ to $\CtrlDMax$.
The appendix reports two quantities on that one sweep, $\sigEP$ and
$\sigma$, which is why both are in view here; the statement that bears on
this section has to be the one about $\sigma$.  They are separate
quantities, each computed from the evaluator's own conditionals, and what
makes each of them non-zero is the same incompatibility.  The behaviour of
one does not follow from the behaviour of the other, and the sweep does not
show them moving together: ranked against each other they give
$\rho = \CtrlRhoXS$ at $p = \CtrlPXS$, which establishes nothing either way
at $n = \CtrlNWidths$.  What $\sigma$ does over the sweep is
not approach zero: $\langle|\sigma|\rangle$ runs between $\CtrlSigAbsMin$
and $\CtrlSigAbsMax$ across all $\CtrlNWidths$ widths.  By the
standard of Section~\ref{sec:window} that is as far as it goes --- a sweep
of $\CtrlNWidths$ points cannot establish that a quantity is constant, and
we do not claim the incompatibility is permanent.  What we claim is that
nothing measured here is approaching the compatible limit.

None of this transfers to $\Delta$.  On a fixed text there is no limit in
which $\Delta$ becomes zero and no identity to appeal to: it is a
log-ratio between a text
and its reversal, and nothing makes that ratio vanish as the autoregressive
evaluator's parameter count grows along the AR ladder.  Measured, its magnitude grows rather than shrinks along the AR ladder
(Section~\ref{sec:capacity}).  The grounds for declining to read $\Delta$ as
a measurement of causal order are the ones given there, in
Section~\ref{sec:resolution} and in Section~\ref{sec:construct}, and they
are empirical: they would stand even if everything this section says about
$\sigma$ were wrong.

The within-text reading and the unrearranged reading are computed by
separate programs, \fn{block_sigma_fb.py} and
\fn{order_intervention.py}.  They are not independent implementations: the
first imports the text loader and the masked-model provider from the second,
so what they share is the input path and the model interface, and what
differs is the construction of $\sigma$ itself.  On the authored
arrangement the median $\sigma$ they return agrees to the printed precision
in all \FbNCorpora\ corpora, which is the check that the rank construction
measures what the direct one does (\fn{check_two_constructions.py}).  That
check covers \texttt{LLaDA} only: no direct run exists for
\texttt{bert-base}.

\subsection{The mask rate \texttt{bert-base} was trained at}
\label{sec:maskrate}

One objection applies to \texttt{bert-base} and not to \texttt{LLaDA}.
Equation~\eqref{eq:aoar} sweeps the fraction of masked positions from $1$
down to $1/L$ as the chain rule reveals tokens.  LLaDA is trained with the
masking fraction drawn uniformly, so every term is in distribution.
\textsc{bert} is trained at a single rate, $0.15$, so most terms of
Eq.~\eqref{eq:aoar} ask it for a conditional at a rate it never saw.  A
natural worry is that the chain rule simply collapses there, and that
anything computed from it is noise.

It does not collapse.  Across all \FbNTexts\ texts and both masked evaluators ---
\FbNScorings\ scorings --- \FbDegenTotal\ produced $n!$ values that were
all equal.  That is what the column detects and all that it detects: a text
whose arrangements the masked evaluator cannot separate at all would have no
midrank, and none occurred, at any mask rate asked for.  It is not a claim
that the $n!$ values are distinct from one another, which this test does not
examine.  The reading is also not degenerate in the weaker
sense: within texts \texttt{bert-base}'s causally labelled group
sits below the null in $\FbCausalBelowBertBase$ of
$\FbNCorporaOfBertBase$ corpora, as \texttt{LLaDA}'s does in
$\FbCausalBelowLLaDA$ of $\FbNCorporaOfLLaDA$.  What differs between the two
masked evaluators within texts is the level the groups sit at, not the order
they sit in.

How far the authored arrangement sits from the null under $\sigma$ is the
quantity Section~\ref{sec:fbresult} measured, and what it found was that
its value tracks the masked evaluator and not the label.  We do not claim
\textsc{bert}'s out-of-distribution conditionals are accurate.  We claim only
that they are usable, that the collapse the objection anticipates does not
occur, and that the conclusion of Section~\ref{sec:fbdef} rests on an
identity rather than on either masked evaluator.
 after Section "Results II".

\section{Results III: $\sigma$, on masked evaluators}
\label{sec:truesigma}

Everything so far measures $\Delta$: a text scored against a rearrangement
of itself, which is the fixed-text reading of
Ref.~\cite{sagawa2026}~\cite[Sec.~V~B]{sagawa2026}.
Section~\ref{sec:notareversal} gave the reason it has no origin on the text
it is read from.  This section measures a comparison that has one.

\subsection{The estimator, and the identity that constrains it}
\label{sec:fbdef}

The estimator is Eq.~\eqref{eq:fb}: the any-order chain rule of
Eq.~\eqref{eq:aoar} evaluated at $\pi = $ left-to-right and again at $\pi
= $ right-to-left on the same text.  Equation~\eqref{eq:aoar} is a genuine
log-probability --- it normalises to one over sequences, for any fixed
$\pi$ --- and both of its evaluations score the same $x$.  What differs is the direction in which the
model assembles it: forward, one token at a time from the left, against
backward, one token at a time from the right.  That is the forward path
against its reverse, which is what the stochastic-thermodynamic definition
asks for.

Equation~\eqref{eq:fb} carries an identity that decides how it can be read,
stated as Eq.~\eqref{eq:fbzero} and proved here.  If the model's
conditionals are those of some joint distribution $P$, then every term of
Eq.~\eqref{eq:aoar} telescopes and $\log P_\pi(x) = \log P(x)$ for every
$\pi$; the two terms of Eq.~\eqref{eq:fb} are then equal, and $\sigma(x) =
0$ for every $x$, causal or not.  The size of a non-zero $\sigma$ is therefore
not a property of the text.  Where one arrangement of a text falls among
the rearrangements of that same text is a separate question, and
Section~\ref{sec:fbrank} is where it is asked.  Whatever
non-zero reading appears is a measure of the masked evaluator's \emph{in}compatibility
--- the failure of its conditionals to derive from any joint --- as exposed
by that text.  This is not an empirical finding to be checked; it is what
Eq.~\eqref{eq:aoar} says.

Two consequences follow immediately.

An autoregressive evaluator has nothing to report here, and it is worth
being exact about why, because the two available answers are different.  The
backward term of Eq.~\eqref{eq:fb} is $\log P_{\mathrm{r2l}}(x)$, and each of
its factors
conditions on a suffix while the prefix is still unknown, and a model that
attends only to what precedes emits none of them.  Equation~\eqref{eq:fb}
is well defined; what an autoregressive evaluator cannot do is supply the
conditionals it asks for.  So $\sigma$ is not evaluable on the
autoregressive evaluators Ref.~\cite{sagawa2026} uses.  What is absent is
not a number but an independent number: a masked model's backward term is
new information because its conditionals derive from no joint, and an
autoregressive model's conditionals derive from one.  That explains the
substitution of $\Delta$ for $\sigma$; it does not license reading $\Delta$
as $\sigma$.

A masked evaluator can be placed here, and is the only kind that can.  It
supplies every conditional Eq.~\eqref{eq:aoar} asks for, in any order.  That
is also the scope of Eq.~\eqref{eq:fbzero}: it says that $\sigma$ vanishes
for an evaluator whose conditionals derive from one joint, and an evaluator
that cannot supply those conditionals is not a case of it either way.  For a
masked evaluator no \emph{one} joint need account for all of the unmasking
orders --- that is what incompatibility is --- and where none does,
Eq.~\eqref{eq:fbzero} does not force $\sigma$ to zero.
Appendix~\ref{sec:toy} measured that incompatibility directly, as the
entropy production $\sigEP$ on a system small enough to enumerate, and
Table~\ref{tab:bayesroute} reads $\sigma$ itself on those same models:
$\sigEP$ and $\sigma$ are both non-zero at every width of that sweep, and
nothing in it relates their sizes (Section~\ref{sec:fbresult} gives the rank
correlation and the $n$ it rests on).  Here $\sigma$ is read off real text.
What the identity fixes is what a non-zero reading means --- that the
conditionals behind it derive from no single joint --- and not what the
value moves with.  Whether it also happens to separate the two label groups
is therefore not settled in advance, and has to be measured.

\subsection{Where the authored arrangement sits}
\label{sec:fbrank}

A single non-zero $\sigma(x)$ says that the conditionals behind it derive
from no single joint, and nothing more; it is not an amount of anything.
The question this section asks is narrower: whether the value $\sigma$
returns is nevertheless higher on the arrangement a text came in than on its
rearrangements, and for that the text has to be compared with itself.  Taking the same block arrangements as
Section~\ref{sec:blkord} --- the $n!$ orders of the same $n$ sentences --- and
writing $x_\pi$ for the text with its blocks in order $\pi$, we score
\begin{equation}
f(\pi) \;=\; \sigma(x_\pi)
\label{eq:fbperm}
\end{equation}
for every $\pi$, and record the midrank of $f(\mathrm{id})$, the authored arrangement,
among the $n!$ values.  Length, vocabulary, topic and style are identical
across the arrangements by construction, so the between-set confounds of
Section~\ref{sec:construct} cannot act.

The null is $0.50$, but only one of the two stages of Section~\ref{sec:blkord} carries over.  There the permutation set itself had
a centre, fixed at zero by the antisymmetry $\Delta(\pi) = -\Delta(\pi \circ
\mathrm{rev})$.  Equation~\eqref{eq:fbperm} has no such symmetry: reversing the arrangement changes which text is scored and does not negate $\sigma$,
so the set of $f(\pi)$ is not centred and supplies no origin of its own.  What survives is the second stage.  Under the hypothesis that the authored arrangement
is one arrangement among $n!$, its rank is uniform, so the midrank has
expectation $0.50$ whatever the shape of the $f$ distribution.

Two situations make the $n!$ values coincide, and both must be reported at
the null rather than ranked.  A masked evaluator whose conditionals derive
from one joint gives $f(\pi) = 0$ for every $\pi$ by
Eq.~\eqref{eq:fbzero}; in floating point those values differ by
$\sim\!10^{-16}$, and ranking them would manufacture a midrank out of
rounding noise.  A text whose sentences are identical strings has
arrangements that are literally the same token sequence.  Values within
$10^{-12}$ are therefore treated as tied, that text is flagged degenerate,
and its midrank is set to $0.50$ exactly.  Table~\ref{tab:fbrank} reports how
many scorings that was: \FbDegenTotal\ of \FbNScorings, over
\FbNTexts\ texts read by each of the two masked evaluators used here,
\texttt{bert-base-uncased} and \texttt{LLaDA-8B-Base}.

This is the \emph{within-text reading}: one text, its $n!$ arrangements
scored by Eq.~\eqref{eq:fbperm}, and the authored arrangement's midrank
among them.  Section~\ref{sec:fbresult} also reads $\sigma$ a second way,
without rearranging anything, and keeps the two apart throughout.


\subsection{What the measurement says}
\label{sec:fbresult}

\begin{table}[tbp]
\centering
{\small% generated by make_tables.py -- do not edit
\begin{tabular}{llrrrrrrrr}
\toprule
evaluator & corpus & $n$ & degen & causal & non-causal & $p_{c}$ & $p_{nc}$ & $r$ & $p$ \\
\midrule
bert-base-uncased & Opus46 & 200 & 0 & $0.458$ & $0.521$ & $0.366$ & $0.968$ & $-0.060$ & $0.462$ \\
bert-base-uncased & GPT54 & 200 & 0 & $0.479$ & $0.521$ & $0.676$ & $0.514$ & $-0.061$ & $0.454$ \\
bert-base-uncased & GPT56 & 200 & 0 & $0.479$ & $0.583$ & $0.989$ & $0.037$ & $-0.126$ & $0.122$ \\
bert-base-uncased & Gemini31 & 200 & 0 & $0.292$ & $0.479$ & $6\times10^{-6}$ & $0.759$ & $-0.287$ & $4\times10^{-4}$ \\
bert-base-uncased & Fable5 & 200 & 0 & $0.417$ & $0.521$ & $0.050$ & $0.736$ & $-0.130$ & $0.112$ \\
LLaDA-8B-Base & Opus46 & 200 & 0 & $0.292$ & $0.229$ & $6\times10^{-6}$ & $1\times10^{-11}$ & $+0.107$ & $0.192$ \\
LLaDA-8B-Base & GPT54 & 200 & 0 & $0.167$ & $0.188$ & $3\times10^{-15}$ & $9\times10^{-14}$ & $-0.187$ & $0.022$ \\
LLaDA-8B-Base & GPT56 & 200 & 0 & $0.229$ & $0.271$ & $2\times10^{-11}$ & $4\times10^{-9}$ & $-0.111$ & $0.175$ \\
LLaDA-8B-Base & Gemini31 & 200 & 0 & $0.146$ & $0.188$ & $9\times10^{-15}$ & $1\times10^{-13}$ & $-0.128$ & $0.115$ \\
LLaDA-8B-Base & Fable5 & 200 & 0 & $0.208$ & $0.229$ & $7\times10^{-11}$ & $6\times10^{-12}$ & $-0.047$ & $0.563$ \\
\bottomrule
\end{tabular}
}
\caption{The forward/backward entropy production of
Eq.~\eqref{eq:fb}, read as the midrank of the authored arrangement among
the $n!$ arrangements of the same text's own sentences.  Columns from the
left: the masked evaluator; the corpus; $n$, the texts in that cell; how
many were degenerate (all $n!$ values tied, hence reported at the null and
excluded from $r$); the median midrank of the causal group and of the
non-causal group; $p_{c}$ and $p_{nc}$, the two-sided Wilcoxon signed-rank
$p$ of each group against the $0.50$ null; $r$, the rank-biserial
correlation of Eq.~\eqref{eq:rb} between the two groups, negative when the
causal group ranks lower; and its two-sided Mann--Whitney $p$.  The two
signed-rank columns answer a different question from the last one: whether
each group is displaced from the null at all, as against whether the groups
differ from each other.  Masked evaluators only
(Section~\ref{sec:truesigma}).}
\label{tab:fbrank}
\end{table}

Table~\ref{tab:fbrank} gives, for each masked evaluator and corpus, the median
midrank of the authored arrangement in each label group, the two-sided
Wilcoxon signed-rank $p$ of each group against the $0.50$ null in the
$p_{c}$ and $p_{nc}$ columns, and the effect size between the groups.  Two
readings come out of it, and they are of different kinds: where the authored
arrangement sits, which turns out to be a property of the masked evaluator
rather than of the text, and what the causal label adds to that, which turns
out to be almost nothing.

\begin{table}[tbp]
\centering
% generated by make_tables.py -- do not edit
\begin{tabular}{lrrrrr}
\toprule
evaluator & authored & steps & causal & steps & $r$ \\
\midrule
bert-base-uncased & $+0.0208$ & $+1.0$ & $-0.1042$ & $-5.0$ & $-0.126$ \\
LLaDA-8B-Base & $-0.2708$ & $-13.0$ & $-0.0208$ & $-1.0$ & $-0.111$ \\
\bottomrule
\end{tabular}

\caption{The same measurement as Table~\ref{tab:fbrank}, summarised per
masked evaluator.  Columns from the left: the masked evaluator; the authored
component $\mathrm{med}(\mathrm{pct}_{\text{non-causal}}) - 0.50$, as a
median over the \NCorp\ corpora; that component in steps of the
$1/\GridK = \GridStep$ grid; the causal component of Eq.~\eqref{eq:caus},
likewise a median over corpora; that component in grid steps; and $r$, the
median over the \NCorp\ corpora of the rank-biserial correlation
Table~\ref{tab:fbrank} prints cell by cell, negative when the causal group
ranks lower.  These are the
numbers the text quotes; Table~\ref{tab:fbrank} is where they come from, one
corpus at a time.}
\label{tab:fbdecomp}
\end{table}

\paragraph{The authored arrangement's midrank is a property of the
masked evaluator.}
Equation~\eqref{eq:auth} is defined on midranks and says nothing about
which quantity produced them, so it reads under $\sigma$ exactly as it does
under $\Delta$.  Under \texttt{LLaDA} the authored component so read is
$\FbAuthoredLLaDA$: the authored arrangement sits that
far below the null, or $\FbAuthoredStepsLLaDA$ steps of the
$1/\GridK = \GridStep$ grid (Table~\ref{tab:fbdecomp}).  It is far from the null in both
label groups: \FbDispNRejectingLLaDA\ of its \FbNTestsDispLLaDA\
distance-from-null tests clear a Bonferroni threshold of
$\FbDispBonferroni$, set for the $\FbNTestsDisp$ such tests this section
reports, the largest of them at $p = \FbDispPMaxLLaDA$.  Those tests are
two-sided, so they establish that it is not at the null and not which side of it; the
direction is the sign of $\FbAuthoredLLaDA$ itself.  Under
\texttt{bert-base} the same texts give $\FbAuthoredBertBase$, which is
$\FbAuthoredStepsBertBase$ steps, and its distance-from-null tests clear that
threshold in $\FbDispNRejectingBertBase$ of $\FbNTestsDispBertBase$ cases.

The two readings stand $\FbStepGap$ grid steps apart on identical input.
Nothing about the texts changed between them; only the masked evaluator
changed.  A
quantity whose reading moves that far when the masked evaluator is replaced
is reporting a property of that evaluator.

The two masked evaluators do agree on one thing, and it is worth separating from
the disagreement.  The causally labelled group sits below the null in
$\FbCausalBelowBertBase$ of $\FbNCorporaOfBertBase$ corpora under
\texttt{bert-base} and in $\FbCausalBelowLLaDA$ of $\FbNCorporaOfLLaDA$
under \texttt{LLaDA}: where the causal group sits \emph{relative to the
null} is the same under both.  That is not a statement about the two groups'
order, which does reverse in one cell: \texttt{LLaDA} on \textsc{Opus46},
the only cell of Table~\ref{tab:fbrank} where $r$ is positive
($r = \FbRMax$).  The
\emph{level} both sit at is not, and the same two counts taken on the
non-causal group show it --- that group is above the null in
$\FbAuthoredAboveBertBase$ corpora under \texttt{bert-base} and in
$\FbAuthoredAboveLLaDA$ under \texttt{LLaDA}.  It is the level that moves by
$\FbStepGap$ steps.

\paragraph{The causal label does essentially nothing.}
Of the \FbNTestsCausal\ masked-evaluator--corpus tests of causal against non-causal,
\FbNSurviving\ survives a Bonferroni threshold of \FbBonferroni;
\FbNNominal\ fall below $0.05$ on their own, against \FbNExpected\ expected
by chance.  The sign is nonetheless the same in \FbNNegative\ of
\FbNCells\ cells, and we report it rather than only the failures: within
each masked evaluator, \texttt{LLaDA} gives \FbSignNegLLaDA\ negative ($p =
\FbSignPLLaDA$, exact binomial) and \texttt{bert-base}
\FbSignNegBertBase\ negative ($p = \FbSignPBertBase$).  Neither reaches
$0.05$ and neither could.  Under the null each corpus is an independent coin,
so the most extreme outcome --- all $\FbNCorporaOfLLaDA$ on one side --- has
probability $1/2^{\FbNCorporaOfLLaDA}$, and the two-sided $p$ of that
outcome is $2/2^{\FbNCorporaOfLLaDA} = \FbSignPBertBase$: the smallest the
test can return at this $n$, and the value \texttt{bert-base} takes.  Pooling the
\FbNCells\ cells would be the wrong test: both masked evaluators score the
\emph{same} texts, so the cells are not independent, and the corpus is the
unit that is.  The one surviving cell is \texttt{bert-base} on
\textsc{Gemini31} --- which is also the corpus where \texttt{LLaDA}'s distance from the null is largest, so it is more plausibly a property of
that corpus than of the label.

Taken together: the authored component under $\sigma$ moves
$\FbStepGap$ grid steps when the masked evaluator is changed, while
exchanging the causal and non-causal labels leaves \FbNSurviving\ of
\FbNCells\ cells surviving correction.  The separation between causally and
arbitrarily ordered texts that $\Delta$ shows in \ContrastNRejecting\ of
\ContrastNCells\ cells does not reappear here.  This is
Eq.~\eqref{eq:fbzero} seen from the measurement side.  The identity says
that where the conditionals derive from a joint there is no value range at
all; what the measurement adds is that where they do not --- which is the
case for \texttt{bert-base} and \texttt{LLaDA} alike --- what moves inside
the range that opens is the evaluator and not the label.

\begin{table}[tbp]
\centering
% generated by make_tables.py -- do not edit
\begin{tabular}{lrr}
\toprule
corpus & bert-base-uncased & LLaDA-8B-Base \\
\midrule
Opus46 & $-0.1235$ & $+0.0574$ \\
GPT54 & $-0.0169$ & $+0.0753$ \\
GPT56 & $+0.0864$ & $+0.1088$ \\
Gemini31 & $-0.1191$ & $+0.0161$ \\
Fable5 & $-0.0384$ & $+0.0552$ \\
\midrule
\multicolumn{3}{l}{\emph{causal component}} \\
Opus46 & $-0.2712$ & $+0.0127$ \\
GPT54 & $-0.4284$ & $-0.0295$ \\
GPT56 & $-0.3251$ & $-0.0500$ \\
Gemini31 & $-0.1727$ & $-0.0284$ \\
Fable5 & $-0.3798$ & $-0.0552$ \\
\bottomrule
\end{tabular}

\caption{The unrearranged reading of Eq.~\eqref{eq:fb}: each text scored
once in each direction, in the arrangement it came in, with no permutation
set and no midrank.  \emph{Upper block}: the median $\sigma$ over all of a
cell's texts, both label groups together, per token.  It is not the authored
component of Eq.~\eqref{eq:auth}, which is a midrank against a null of
$0.50$; the null here is $\sigma = 0$.  \emph{Lower block}: the causal
component of this reading --- the median $\sigma$ of the causally labelled
texts less that of the non-causal ones --- which shares its name with
Eq.~\eqref{eq:caus} but is a difference of raw $\sigma$ values rather than
of midranks.  For a masked evaluator whose conditionals derive from a joint
every entry of both blocks would be exactly zero by Eq.~\eqref{eq:fbzero}.
None is, and on \FbBetweenNOpposite\ of the \FbBetweenNCompared\ corpora
the two masked evaluators give the upper block's median $\sigma$ opposite
signs.}
\label{tab:fbbetween}
\end{table}

A second reading makes the same point more directly.  Call it the
\emph{unrearranged reading}: Eq.~\eqref{eq:fb} evaluated as it stands, on
the text in the arrangement it came in, with no permutation set behind it.
Across cells the median $\sigma$
runs $\FbBetweenMinLLaDA$ to $\FbBetweenMaxLLaDA$ per token under
\texttt{LLaDA} and $\FbBetweenMinBertBase$ to $\FbBetweenMaxBertBase$ under
\texttt{bert-base}, positive in \FbBetweenNPosLLaDA\ of
\FbBetweenNOfLLaDA\ corpora under the first and \FbBetweenNPosBertBase\ of
\FbBetweenNOfBertBase\ under the second (Table~\ref{tab:fbbetween}).  On
\FbBetweenNOpposite\ of the \FbBetweenNCompared\ corpora the two give that
median opposite signs.  The texts are identical; what differs is which masked
evaluator computed $\sigma$.  What Eq.~\eqref{eq:fbzero} settles before any
corpus is scored is that a non-zero reading belongs to the evaluator's
incompatibility rather than to the text; a disagreement in sign on identical
texts is one way that shows, and confirms it rather than establishing it.  The causal
component of the same reading is the lower block of
Table~\ref{tab:fbbetween}: a median of $\FbBetweenCausalLLaDA$ under
\texttt{LLaDA} and $\FbBetweenCausalBertBase$ under \texttt{bert-base},
negative in \FbBetweenNNegOfLLaDA\ and \FbBetweenNNegOfBertBase\ corpora
respectively.  Both readings therefore give the causal group's median less
the non-causal group's the same sign, as a median over corpora: the causally
labelled texts sit lower.  The two are not in the same units: the
within-text reading is a difference of midranks ($\FbCausalLLaDA$ and
$\FbCausalBertBase$, or $\FbCausalStepsLLaDA$ and
$\FbCausalStepsBertBase$ grid steps), the unrearranged reading a difference
of raw $\sigma$.  The agreement is therefore one of direction and not of
size, and neither reaches significance.

\paragraph{Why this is not the argument made against $\Delta$.}
The two negative results in this paper are not of the same kind, and the
distinction decides what either of them licenses.  Under $\Delta$
(Sections~\ref{sec:authored} and~\ref{sec:causal}) the two label groups do
separate, in \ContrastNRejecting\ of \ContrastNCells\ cells.  What cannot be
read off is the size of the part attributable to the causal label: it sits at
the grid's resolution, and the groups have no criterion distinguishing them
beyond the prompts that produced them.  Under $\sigma$
(this section) what a non-zero reading reports is the masked evaluator's
incompatibility, by the identity, and not causal order.  The $\Delta$
result is empirical; the $\sigma$ one follows from the identity.  What $\sigma$ reports is
a defect of the masked evaluator: Eq.~\eqref{eq:fbzero} makes it identically zero
wherever the conditionals derive from a joint, so a reading of any size is
a reading of incompatibility, and a contrast between text sets under
$\sigma$ would be a contrast in how much of that incompatibility each set
exposes.  This is an argument about the quantity and it holds before any
measurement.  The direction differs between the two: under $\Delta$ the
causal group scores higher, under $\sigma$ lower.  We do not account for
that difference, and no claim in this paper needs to --- by Eq.~\eqref{eq:fbzero} a
reading of $\sigma$ of any size, in either direction, is a reading of the
masked evaluator's incompatibility and not of the text's causal structure.

Nor does the sweep of Appendix~\ref{sec:toy} show $\sigma$ heading towards
zero as the transformer width $d$ grows from $\CtrlDMin$ to $\CtrlDMax$.
The appendix reports two quantities on that one sweep, $\sigEP$ and
$\sigma$, which is why both are in view here; the statement that bears on
this section has to be the one about $\sigma$.  They are separate
quantities, each computed from the evaluator's own conditionals, and what
makes each of them non-zero is the same incompatibility.  The behaviour of
one does not follow from the behaviour of the other, and the sweep does not
show them moving together: ranked against each other they give
$\rho = \CtrlRhoXS$ at $p = \CtrlPXS$, which establishes nothing either way
at $n = \CtrlNWidths$.  What $\sigma$ does over the sweep is
not approach zero: $\langle|\sigma|\rangle$ runs between $\CtrlSigAbsMin$
and $\CtrlSigAbsMax$ across all $\CtrlNWidths$ widths.  By the
standard of Section~\ref{sec:window} that is as far as it goes --- a sweep
of $\CtrlNWidths$ points cannot establish that a quantity is constant, and
we do not claim the incompatibility is permanent.  What we claim is that
nothing measured here is approaching the compatible limit.

None of this transfers to $\Delta$.  On a fixed text there is no limit in
which $\Delta$ becomes zero and no identity to appeal to: it is a
log-ratio between a text
and its reversal, and nothing makes that ratio vanish as the autoregressive
evaluator's parameter count grows along the AR ladder.  Measured, its magnitude grows rather than shrinks along the AR ladder
(Section~\ref{sec:capacity}).  The grounds for declining to read $\Delta$ as
a measurement of causal order are the ones given there, in
Section~\ref{sec:resolution} and in Section~\ref{sec:construct}, and they
are empirical: they would stand even if everything this section says about
$\sigma$ were wrong.

The within-text reading and the unrearranged reading are computed by
separate programs, \fn{block_sigma_fb.py} and
\fn{order_intervention.py}.  They are not independent implementations: the
first imports the text loader and the masked-model provider from the second,
so what they share is the input path and the model interface, and what
differs is the construction of $\sigma$ itself.  On the authored
arrangement the median $\sigma$ they return agrees to the printed precision
in all \FbNCorpora\ corpora, which is the check that the rank construction
measures what the direct one does (\fn{check_two_constructions.py}).  That
check covers \texttt{LLaDA} only: no direct run exists for
\texttt{bert-base}.

\subsection{The mask rate \texttt{bert-base} was trained at}
\label{sec:maskrate}

One objection applies to \texttt{bert-base} and not to \texttt{LLaDA}.
Equation~\eqref{eq:aoar} sweeps the fraction of masked positions from $1$
down to $1/L$ as the chain rule reveals tokens.  LLaDA is trained with the
masking fraction drawn uniformly, so every term is in distribution.
\textsc{bert} is trained at a single rate, $0.15$, so most terms of
Eq.~\eqref{eq:aoar} ask it for a conditional at a rate it never saw.  A
natural worry is that the chain rule simply collapses there, and that
anything computed from it is noise.

It does not collapse.  Across all \FbNTexts\ texts and both masked evaluators ---
\FbNScorings\ scorings --- \FbDegenTotal\ produced $n!$ values that were
all equal.  That is what the column detects and all that it detects: a text
whose arrangements the masked evaluator cannot separate at all would have no
midrank, and none occurred, at any mask rate asked for.  It is not a claim
that the $n!$ values are distinct from one another, which this test does not
examine.  The reading is also not degenerate in the weaker
sense: within texts \texttt{bert-base}'s causally labelled group
sits below the null in $\FbCausalBelowBertBase$ of
$\FbNCorporaOfBertBase$ corpora, as \texttt{LLaDA}'s does in
$\FbCausalBelowLLaDA$ of $\FbNCorporaOfLLaDA$.  What differs between the two
masked evaluators within texts is the level the groups sit at, not the order
they sit in.

How far the authored arrangement sits from the null under $\sigma$ is the
quantity Section~\ref{sec:fbresult} measured, and what it found was that
its value tracks the masked evaluator and not the label.  We do not claim
\textsc{bert}'s out-of-distribution conditionals are accurate.  We claim only
that they are usable, that the collapse the objection anticipates does not
occur, and that the conclusion of Section~\ref{sec:fbdef} rests on an
identity rather than on either masked evaluator.
 after Section "Results II".

\section{Results III: $\sigma$, on masked evaluators}
\label{sec:truesigma}

Everything so far measures $\Delta$: a text scored against a rearrangement
of itself, which is the fixed-text reading of
Ref.~\cite{sagawa2026}~\cite[Sec.~V~B]{sagawa2026}.
Section~\ref{sec:notareversal} gave the reason it has no origin on the text
it is read from.  This section measures a comparison that has one.

\subsection{The estimator, and the identity that constrains it}
\label{sec:fbdef}

The estimator is Eq.~\eqref{eq:fb}: the any-order chain rule of
Eq.~\eqref{eq:aoar} evaluated at $\pi = $ left-to-right and again at $\pi
= $ right-to-left on the same text.  Equation~\eqref{eq:aoar} is a genuine
log-probability --- it normalises to one over sequences, for any fixed
$\pi$ --- and both of its evaluations score the same $x$.  What differs is the direction in which the
model assembles it: forward, one token at a time from the left, against
backward, one token at a time from the right.  That is the forward path
against its reverse, which is what the stochastic-thermodynamic definition
asks for.

Equation~\eqref{eq:fb} carries an identity that decides how it can be read,
stated as Eq.~\eqref{eq:fbzero} and proved here.  If the model's
conditionals are those of some joint distribution $P$, then every term of
Eq.~\eqref{eq:aoar} telescopes and $\log P_\pi(x) = \log P(x)$ for every
$\pi$; the two terms of Eq.~\eqref{eq:fb} are then equal, and $\sigma(x) =
0$ for every $x$, causal or not.  The size of a non-zero $\sigma$ is therefore
not a property of the text.  Where one arrangement of a text falls among
the rearrangements of that same text is a separate question, and
Section~\ref{sec:fbrank} is where it is asked.  Whatever
non-zero reading appears is a measure of the masked evaluator's \emph{in}compatibility
--- the failure of its conditionals to derive from any joint --- as exposed
by that text.  This is not an empirical finding to be checked; it is what
Eq.~\eqref{eq:aoar} says.

Two consequences follow immediately.

An autoregressive evaluator has nothing to report here, and it is worth
being exact about why, because the two available answers are different.  The
backward term of Eq.~\eqref{eq:fb} is $\log P_{\mathrm{r2l}}(x)$, and each of
its factors
conditions on a suffix while the prefix is still unknown, and a model that
attends only to what precedes emits none of them.  Equation~\eqref{eq:fb}
is well defined; what an autoregressive evaluator cannot do is supply the
conditionals it asks for.  So $\sigma$ is not evaluable on the
autoregressive evaluators Ref.~\cite{sagawa2026} uses.  What is absent is
not a number but an independent number: a masked model's backward term is
new information because its conditionals derive from no joint, and an
autoregressive model's conditionals derive from one.  That explains the
substitution of $\Delta$ for $\sigma$; it does not license reading $\Delta$
as $\sigma$.

A masked evaluator can be placed here, and is the only kind that can.  It
supplies every conditional Eq.~\eqref{eq:aoar} asks for, in any order.  That
is also the scope of Eq.~\eqref{eq:fbzero}: it says that $\sigma$ vanishes
for an evaluator whose conditionals derive from one joint, and an evaluator
that cannot supply those conditionals is not a case of it either way.  For a
masked evaluator no \emph{one} joint need account for all of the unmasking
orders --- that is what incompatibility is --- and where none does,
Eq.~\eqref{eq:fbzero} does not force $\sigma$ to zero.
Appendix~\ref{sec:toy} measured that incompatibility directly, as the
entropy production $\sigEP$ on a system small enough to enumerate, and
Table~\ref{tab:bayesroute} reads $\sigma$ itself on those same models:
$\sigEP$ and $\sigma$ are both non-zero at every width of that sweep, and
nothing in it relates their sizes (Section~\ref{sec:fbresult} gives the rank
correlation and the $n$ it rests on).  Here $\sigma$ is read off real text.
What the identity fixes is what a non-zero reading means --- that the
conditionals behind it derive from no single joint --- and not what the
value moves with.  Whether it also happens to separate the two label groups
is therefore not settled in advance, and has to be measured.

\subsection{Where the authored arrangement sits}
\label{sec:fbrank}

A single non-zero $\sigma(x)$ says that the conditionals behind it derive
from no single joint, and nothing more; it is not an amount of anything.
The question this section asks is narrower: whether the value $\sigma$
returns is nevertheless higher on the arrangement a text came in than on its
rearrangements, and for that the text has to be compared with itself.  Taking the same block arrangements as
Section~\ref{sec:blkord} --- the $n!$ orders of the same $n$ sentences --- and
writing $x_\pi$ for the text with its blocks in order $\pi$, we score
\begin{equation}
f(\pi) \;=\; \sigma(x_\pi)
\label{eq:fbperm}
\end{equation}
for every $\pi$, and record the midrank of $f(\mathrm{id})$, the authored arrangement,
among the $n!$ values.  Length, vocabulary, topic and style are identical
across the arrangements by construction, so the between-set confounds of
Section~\ref{sec:construct} cannot act.

The null is $0.50$, but only one of the two stages of Section~\ref{sec:blkord} carries over.  There the permutation set itself had
a centre, fixed at zero by the antisymmetry $\Delta(\pi) = -\Delta(\pi \circ
\mathrm{rev})$.  Equation~\eqref{eq:fbperm} has no such symmetry: reversing the arrangement changes which text is scored and does not negate $\sigma$,
so the set of $f(\pi)$ is not centred and supplies no origin of its own.  What survives is the second stage.  Under the hypothesis that the authored arrangement
is one arrangement among $n!$, its rank is uniform, so the midrank has
expectation $0.50$ whatever the shape of the $f$ distribution.

Two situations make the $n!$ values coincide, and both must be reported at
the null rather than ranked.  A masked evaluator whose conditionals derive
from one joint gives $f(\pi) = 0$ for every $\pi$ by
Eq.~\eqref{eq:fbzero}; in floating point those values differ by
$\sim\!10^{-16}$, and ranking them would manufacture a midrank out of
rounding noise.  A text whose sentences are identical strings has
arrangements that are literally the same token sequence.  Values within
$10^{-12}$ are therefore treated as tied, that text is flagged degenerate,
and its midrank is set to $0.50$ exactly.  Table~\ref{tab:fbrank} reports how
many scorings that was: \FbDegenTotal\ of \FbNScorings, over
\FbNTexts\ texts read by each of the two masked evaluators used here,
\texttt{bert-base-uncased} and \texttt{LLaDA-8B-Base}.

This is the \emph{within-text reading}: one text, its $n!$ arrangements
scored by Eq.~\eqref{eq:fbperm}, and the authored arrangement's midrank
among them.  Section~\ref{sec:fbresult} also reads $\sigma$ a second way,
without rearranging anything, and keeps the two apart throughout.


\subsection{What the measurement says}
\label{sec:fbresult}

\begin{table}[tbp]
\centering
{\small% generated by make_tables.py -- do not edit
\begin{tabular}{llrrrrrrrr}
\toprule
evaluator & corpus & $n$ & degen & causal & non-causal & $p_{c}$ & $p_{nc}$ & $r$ & $p$ \\
\midrule
bert-base-uncased & Opus46 & 200 & 0 & $0.458$ & $0.521$ & $0.366$ & $0.968$ & $-0.060$ & $0.462$ \\
bert-base-uncased & GPT54 & 200 & 0 & $0.479$ & $0.521$ & $0.676$ & $0.514$ & $-0.061$ & $0.454$ \\
bert-base-uncased & GPT56 & 200 & 0 & $0.479$ & $0.583$ & $0.989$ & $0.037$ & $-0.126$ & $0.122$ \\
bert-base-uncased & Gemini31 & 200 & 0 & $0.292$ & $0.479$ & $6\times10^{-6}$ & $0.759$ & $-0.287$ & $4\times10^{-4}$ \\
bert-base-uncased & Fable5 & 200 & 0 & $0.417$ & $0.521$ & $0.050$ & $0.736$ & $-0.130$ & $0.112$ \\
LLaDA-8B-Base & Opus46 & 200 & 0 & $0.292$ & $0.229$ & $6\times10^{-6}$ & $1\times10^{-11}$ & $+0.107$ & $0.192$ \\
LLaDA-8B-Base & GPT54 & 200 & 0 & $0.167$ & $0.188$ & $3\times10^{-15}$ & $9\times10^{-14}$ & $-0.187$ & $0.022$ \\
LLaDA-8B-Base & GPT56 & 200 & 0 & $0.229$ & $0.271$ & $2\times10^{-11}$ & $4\times10^{-9}$ & $-0.111$ & $0.175$ \\
LLaDA-8B-Base & Gemini31 & 200 & 0 & $0.146$ & $0.188$ & $9\times10^{-15}$ & $1\times10^{-13}$ & $-0.128$ & $0.115$ \\
LLaDA-8B-Base & Fable5 & 200 & 0 & $0.208$ & $0.229$ & $7\times10^{-11}$ & $6\times10^{-12}$ & $-0.047$ & $0.563$ \\
\bottomrule
\end{tabular}
}
\caption{The forward/backward entropy production of
Eq.~\eqref{eq:fb}, read as the midrank of the authored arrangement among
the $n!$ arrangements of the same text's own sentences.  Columns from the
left: the masked evaluator; the corpus; $n$, the texts in that cell; how
many were degenerate (all $n!$ values tied, hence reported at the null and
excluded from $r$); the median midrank of the causal group and of the
non-causal group; $p_{c}$ and $p_{nc}$, the two-sided Wilcoxon signed-rank
$p$ of each group against the $0.50$ null; $r$, the rank-biserial
correlation of Eq.~\eqref{eq:rb} between the two groups, negative when the
causal group ranks lower; and its two-sided Mann--Whitney $p$.  The two
signed-rank columns answer a different question from the last one: whether
each group is displaced from the null at all, as against whether the groups
differ from each other.  Masked evaluators only
(Section~\ref{sec:truesigma}).}
\label{tab:fbrank}
\end{table}

Table~\ref{tab:fbrank} gives, for each masked evaluator and corpus, the median
midrank of the authored arrangement in each label group, the two-sided
Wilcoxon signed-rank $p$ of each group against the $0.50$ null in the
$p_{c}$ and $p_{nc}$ columns, and the effect size between the groups.  Two
readings come out of it, and they are of different kinds: where the authored
arrangement sits, which turns out to be a property of the masked evaluator
rather than of the text, and what the causal label adds to that, which turns
out to be almost nothing.

\begin{table}[tbp]
\centering
% generated by make_tables.py -- do not edit
\begin{tabular}{lrrrrr}
\toprule
evaluator & authored & steps & causal & steps & $r$ \\
\midrule
bert-base-uncased & $+0.0208$ & $+1.0$ & $-0.1042$ & $-5.0$ & $-0.126$ \\
LLaDA-8B-Base & $-0.2708$ & $-13.0$ & $-0.0208$ & $-1.0$ & $-0.111$ \\
\bottomrule
\end{tabular}

\caption{The same measurement as Table~\ref{tab:fbrank}, summarised per
masked evaluator.  Columns from the left: the masked evaluator; the authored
component $\mathrm{med}(\mathrm{pct}_{\text{non-causal}}) - 0.50$, as a
median over the \NCorp\ corpora; that component in steps of the
$1/\GridK = \GridStep$ grid; the causal component of Eq.~\eqref{eq:caus},
likewise a median over corpora; that component in grid steps; and $r$, the
median over the \NCorp\ corpora of the rank-biserial correlation
Table~\ref{tab:fbrank} prints cell by cell, negative when the causal group
ranks lower.  These are the
numbers the text quotes; Table~\ref{tab:fbrank} is where they come from, one
corpus at a time.}
\label{tab:fbdecomp}
\end{table}

\paragraph{The authored arrangement's midrank is a property of the
masked evaluator.}
Equation~\eqref{eq:auth} is defined on midranks and says nothing about
which quantity produced them, so it reads under $\sigma$ exactly as it does
under $\Delta$.  Under \texttt{LLaDA} the authored component so read is
$\FbAuthoredLLaDA$: the authored arrangement sits that
far below the null, or $\FbAuthoredStepsLLaDA$ steps of the
$1/\GridK = \GridStep$ grid (Table~\ref{tab:fbdecomp}).  It is far from the null in both
label groups: \FbDispNRejectingLLaDA\ of its \FbNTestsDispLLaDA\
distance-from-null tests clear a Bonferroni threshold of
$\FbDispBonferroni$, set for the $\FbNTestsDisp$ such tests this section
reports, the largest of them at $p = \FbDispPMaxLLaDA$.  Those tests are
two-sided, so they establish that it is not at the null and not which side of it; the
direction is the sign of $\FbAuthoredLLaDA$ itself.  Under
\texttt{bert-base} the same texts give $\FbAuthoredBertBase$, which is
$\FbAuthoredStepsBertBase$ steps, and its distance-from-null tests clear that
threshold in $\FbDispNRejectingBertBase$ of $\FbNTestsDispBertBase$ cases.

The two readings stand $\FbStepGap$ grid steps apart on identical input.
Nothing about the texts changed between them; only the masked evaluator
changed.  A
quantity whose reading moves that far when the masked evaluator is replaced
is reporting a property of that evaluator.

The two masked evaluators do agree on one thing, and it is worth separating from
the disagreement.  The causally labelled group sits below the null in
$\FbCausalBelowBertBase$ of $\FbNCorporaOfBertBase$ corpora under
\texttt{bert-base} and in $\FbCausalBelowLLaDA$ of $\FbNCorporaOfLLaDA$
under \texttt{LLaDA}: where the causal group sits \emph{relative to the
null} is the same under both.  That is not a statement about the two groups'
order, which does reverse in one cell: \texttt{LLaDA} on \textsc{Opus46},
the only cell of Table~\ref{tab:fbrank} where $r$ is positive
($r = \FbRMax$).  The
\emph{level} both sit at is not, and the same two counts taken on the
non-causal group show it --- that group is above the null in
$\FbAuthoredAboveBertBase$ corpora under \texttt{bert-base} and in
$\FbAuthoredAboveLLaDA$ under \texttt{LLaDA}.  It is the level that moves by
$\FbStepGap$ steps.

\paragraph{The causal label does essentially nothing.}
Of the \FbNTestsCausal\ masked-evaluator--corpus tests of causal against non-causal,
\FbNSurviving\ survives a Bonferroni threshold of \FbBonferroni;
\FbNNominal\ fall below $0.05$ on their own, against \FbNExpected\ expected
by chance.  The sign is nonetheless the same in \FbNNegative\ of
\FbNCells\ cells, and we report it rather than only the failures: within
each masked evaluator, \texttt{LLaDA} gives \FbSignNegLLaDA\ negative ($p =
\FbSignPLLaDA$, exact binomial) and \texttt{bert-base}
\FbSignNegBertBase\ negative ($p = \FbSignPBertBase$).  Neither reaches
$0.05$ and neither could.  Under the null each corpus is an independent coin,
so the most extreme outcome --- all $\FbNCorporaOfLLaDA$ on one side --- has
probability $1/2^{\FbNCorporaOfLLaDA}$, and the two-sided $p$ of that
outcome is $2/2^{\FbNCorporaOfLLaDA} = \FbSignPBertBase$: the smallest the
test can return at this $n$, and the value \texttt{bert-base} takes.  Pooling the
\FbNCells\ cells would be the wrong test: both masked evaluators score the
\emph{same} texts, so the cells are not independent, and the corpus is the
unit that is.  The one surviving cell is \texttt{bert-base} on
\textsc{Gemini31} --- which is also the corpus where \texttt{LLaDA}'s distance from the null is largest, so it is more plausibly a property of
that corpus than of the label.

Taken together: the authored component under $\sigma$ moves
$\FbStepGap$ grid steps when the masked evaluator is changed, while
exchanging the causal and non-causal labels leaves \FbNSurviving\ of
\FbNCells\ cells surviving correction.  The separation between causally and
arbitrarily ordered texts that $\Delta$ shows in \ContrastNRejecting\ of
\ContrastNCells\ cells does not reappear here.  This is
Eq.~\eqref{eq:fbzero} seen from the measurement side.  The identity says
that where the conditionals derive from a joint there is no value range at
all; what the measurement adds is that where they do not --- which is the
case for \texttt{bert-base} and \texttt{LLaDA} alike --- what moves inside
the range that opens is the evaluator and not the label.

\begin{table}[tbp]
\centering
% generated by make_tables.py -- do not edit
\begin{tabular}{lrr}
\toprule
corpus & bert-base-uncased & LLaDA-8B-Base \\
\midrule
Opus46 & $-0.1235$ & $+0.0574$ \\
GPT54 & $-0.0169$ & $+0.0753$ \\
GPT56 & $+0.0864$ & $+0.1088$ \\
Gemini31 & $-0.1191$ & $+0.0161$ \\
Fable5 & $-0.0384$ & $+0.0552$ \\
\midrule
\multicolumn{3}{l}{\emph{causal component}} \\
Opus46 & $-0.2712$ & $+0.0127$ \\
GPT54 & $-0.4284$ & $-0.0295$ \\
GPT56 & $-0.3251$ & $-0.0500$ \\
Gemini31 & $-0.1727$ & $-0.0284$ \\
Fable5 & $-0.3798$ & $-0.0552$ \\
\bottomrule
\end{tabular}

\caption{The unrearranged reading of Eq.~\eqref{eq:fb}: each text scored
once in each direction, in the arrangement it came in, with no permutation
set and no midrank.  \emph{Upper block}: the median $\sigma$ over all of a
cell's texts, both label groups together, per token.  It is not the authored
component of Eq.~\eqref{eq:auth}, which is a midrank against a null of
$0.50$; the null here is $\sigma = 0$.  \emph{Lower block}: the causal
component of this reading --- the median $\sigma$ of the causally labelled
texts less that of the non-causal ones --- which shares its name with
Eq.~\eqref{eq:caus} but is a difference of raw $\sigma$ values rather than
of midranks.  For a masked evaluator whose conditionals derive from a joint
every entry of both blocks would be exactly zero by Eq.~\eqref{eq:fbzero}.
None is, and on \FbBetweenNOpposite\ of the \FbBetweenNCompared\ corpora
the two masked evaluators give the upper block's median $\sigma$ opposite
signs.}
\label{tab:fbbetween}
\end{table}

A second reading makes the same point more directly.  Call it the
\emph{unrearranged reading}: Eq.~\eqref{eq:fb} evaluated as it stands, on
the text in the arrangement it came in, with no permutation set behind it.
Across cells the median $\sigma$
runs $\FbBetweenMinLLaDA$ to $\FbBetweenMaxLLaDA$ per token under
\texttt{LLaDA} and $\FbBetweenMinBertBase$ to $\FbBetweenMaxBertBase$ under
\texttt{bert-base}, positive in \FbBetweenNPosLLaDA\ of
\FbBetweenNOfLLaDA\ corpora under the first and \FbBetweenNPosBertBase\ of
\FbBetweenNOfBertBase\ under the second (Table~\ref{tab:fbbetween}).  On
\FbBetweenNOpposite\ of the \FbBetweenNCompared\ corpora the two give that
median opposite signs.  The texts are identical; what differs is which masked
evaluator computed $\sigma$.  What Eq.~\eqref{eq:fbzero} settles before any
corpus is scored is that a non-zero reading belongs to the evaluator's
incompatibility rather than to the text; a disagreement in sign on identical
texts is one way that shows, and confirms it rather than establishing it.  The causal
component of the same reading is the lower block of
Table~\ref{tab:fbbetween}: a median of $\FbBetweenCausalLLaDA$ under
\texttt{LLaDA} and $\FbBetweenCausalBertBase$ under \texttt{bert-base},
negative in \FbBetweenNNegOfLLaDA\ and \FbBetweenNNegOfBertBase\ corpora
respectively.  Both readings therefore give the causal group's median less
the non-causal group's the same sign, as a median over corpora: the causally
labelled texts sit lower.  The two are not in the same units: the
within-text reading is a difference of midranks ($\FbCausalLLaDA$ and
$\FbCausalBertBase$, or $\FbCausalStepsLLaDA$ and
$\FbCausalStepsBertBase$ grid steps), the unrearranged reading a difference
of raw $\sigma$.  The agreement is therefore one of direction and not of
size, and neither reaches significance.

\paragraph{Why this is not the argument made against $\Delta$.}
The two negative results in this paper are not of the same kind, and the
distinction decides what either of them licenses.  Under $\Delta$
(Sections~\ref{sec:authored} and~\ref{sec:causal}) the two label groups do
separate, in \ContrastNRejecting\ of \ContrastNCells\ cells.  What cannot be
read off is the size of the part attributable to the causal label: it sits at
the grid's resolution, and the groups have no criterion distinguishing them
beyond the prompts that produced them.  Under $\sigma$
(this section) what a non-zero reading reports is the masked evaluator's
incompatibility, by the identity, and not causal order.  The $\Delta$
result is empirical; the $\sigma$ one follows from the identity.  What $\sigma$ reports is
a defect of the masked evaluator: Eq.~\eqref{eq:fbzero} makes it identically zero
wherever the conditionals derive from a joint, so a reading of any size is
a reading of incompatibility, and a contrast between text sets under
$\sigma$ would be a contrast in how much of that incompatibility each set
exposes.  This is an argument about the quantity and it holds before any
measurement.  The direction differs between the two: under $\Delta$ the
causal group scores higher, under $\sigma$ lower.  We do not account for
that difference, and no claim in this paper needs to --- by Eq.~\eqref{eq:fbzero} a
reading of $\sigma$ of any size, in either direction, is a reading of the
masked evaluator's incompatibility and not of the text's causal structure.

Nor does the sweep of Appendix~\ref{sec:toy} show $\sigma$ heading towards
zero as the transformer width $d$ grows from $\CtrlDMin$ to $\CtrlDMax$.
The appendix reports two quantities on that one sweep, $\sigEP$ and
$\sigma$, which is why both are in view here; the statement that bears on
this section has to be the one about $\sigma$.  They are separate
quantities, each computed from the evaluator's own conditionals, and what
makes each of them non-zero is the same incompatibility.  The behaviour of
one does not follow from the behaviour of the other, and the sweep does not
show them moving together: ranked against each other they give
$\rho = \CtrlRhoXS$ at $p = \CtrlPXS$, which establishes nothing either way
at $n = \CtrlNWidths$.  What $\sigma$ does over the sweep is
not approach zero: $\langle|\sigma|\rangle$ runs between $\CtrlSigAbsMin$
and $\CtrlSigAbsMax$ across all $\CtrlNWidths$ widths.  By the
standard of Section~\ref{sec:window} that is as far as it goes --- a sweep
of $\CtrlNWidths$ points cannot establish that a quantity is constant, and
we do not claim the incompatibility is permanent.  What we claim is that
nothing measured here is approaching the compatible limit.

None of this transfers to $\Delta$.  On a fixed text there is no limit in
which $\Delta$ becomes zero and no identity to appeal to: it is a
log-ratio between a text
and its reversal, and nothing makes that ratio vanish as the autoregressive
evaluator's parameter count grows along the AR ladder.  Measured, its magnitude grows rather than shrinks along the AR ladder
(Section~\ref{sec:capacity}).  The grounds for declining to read $\Delta$ as
a measurement of causal order are the ones given there, in
Section~\ref{sec:resolution} and in Section~\ref{sec:construct}, and they
are empirical: they would stand even if everything this section says about
$\sigma$ were wrong.

The within-text reading and the unrearranged reading are computed by
separate programs, \fn{block_sigma_fb.py} and
\fn{order_intervention.py}.  They are not independent implementations: the
first imports the text loader and the masked-model provider from the second,
so what they share is the input path and the model interface, and what
differs is the construction of $\sigma$ itself.  On the authored
arrangement the median $\sigma$ they return agrees to the printed precision
in all \FbNCorpora\ corpora, which is the check that the rank construction
measures what the direct one does (\fn{check_two_constructions.py}).  That
check covers \texttt{LLaDA} only: no direct run exists for
\texttt{bert-base}.

\subsection{The mask rate \texttt{bert-base} was trained at}
\label{sec:maskrate}

One objection applies to \texttt{bert-base} and not to \texttt{LLaDA}.
Equation~\eqref{eq:aoar} sweeps the fraction of masked positions from $1$
down to $1/L$ as the chain rule reveals tokens.  LLaDA is trained with the
masking fraction drawn uniformly, so every term is in distribution.
\textsc{bert} is trained at a single rate, $0.15$, so most terms of
Eq.~\eqref{eq:aoar} ask it for a conditional at a rate it never saw.  A
natural worry is that the chain rule simply collapses there, and that
anything computed from it is noise.

It does not collapse.  Across all \FbNTexts\ texts and both masked evaluators ---
\FbNScorings\ scorings --- \FbDegenTotal\ produced $n!$ values that were
all equal.  That is what the column detects and all that it detects: a text
whose arrangements the masked evaluator cannot separate at all would have no
midrank, and none occurred, at any mask rate asked for.  It is not a claim
that the $n!$ values are distinct from one another, which this test does not
examine.  The reading is also not degenerate in the weaker
sense: within texts \texttt{bert-base}'s causally labelled group
sits below the null in $\FbCausalBelowBertBase$ of
$\FbNCorporaOfBertBase$ corpora, as \texttt{LLaDA}'s does in
$\FbCausalBelowLLaDA$ of $\FbNCorporaOfLLaDA$.  What differs between the two
masked evaluators within texts is the level the groups sit at, not the order
they sit in.

How far the authored arrangement sits from the null under $\sigma$ is the
quantity Section~\ref{sec:fbresult} measured, and what it found was that
its value tracks the masked evaluator and not the label.  We do not claim
\textsc{bert}'s out-of-distribution conditionals are accurate.  We claim only
that they are usable, that the collapse the objection anticipates does not
occur, and that the conclusion of Section~\ref{sec:fbdef} rests on an
identity rather than on either masked evaluator.
